\subsection{A Reply to the Papal Encyclical of Pope Leo XIII, on
Reunion}\label{a-reply-to-the-papal-encyclical-of-pope-leo-xiii-on-reunion}

\emph{To the most Sacred and Most Divinely-beloved Brethren in Christ
the Metropolitans and Bishops, and their sacred and venerable Clergy,
and all the godly and orthodox Laity of the Most Holy Apostolic and
Patriarchal Throne of Constantinople.}

\begin{quote}
\emph{"Remember them which have the rule over you, who have spoken unto
you the word of God: whose faith follow, considering the end of their
own conversation:}

\emph{"Jesus Christ the same yesterday, and today, and for ever. Be not
carried about with divers and strange doctrines."} (Heb. xiii. 7, 8).
\end{quote}

I. Every godly and orthodox soul, which has a sincere zeal for the glory
of God, is deeply afflicted and weighed down with great pain upon seeing
that he, who detests that which is good and is a murderer from the
beginning, impelled by envy of man\textquotesingle s salvation, never
ceases continually to sow divers tares in the field of the Lord, in
order to sift the wheat. From this source indeed, even from the earliest
times, there sprang up in the Church of God heretical tares, which have
in many ways made havoc, and do still make havoc, of the salvation of
mankind by Christ; which moreover, as bad seeds and corrupted members,
are rightly cut off from the sound body of the orthodox catholic Church
of Christ. But in these last times the evil one has rent from the
orthodox Church of Christ even whole nations in the West, having
inflated the bishops of Rome with thoughts of excessive arrogance, which
has given birth to divers lawless and anti-evangelical innovations. And
not only so, but furthermore the Popes of Rome from time to time,
pursuing absolutely and without examination modes of union according to
their own fancy, strive by every means to reduce to their own errors the
catholic Church of Christ, which throughout the world walks unshaken in
the orthodoxy of faith transmitted to her by the Fathers.

II. Accordingly the Pope of Rome, Leo XIII, on the occasion of his
episcopal jubilee, published in the month of June of the year of grace
1895 an encyclical letter, addressed to the leaders and peoples of the
world, by which he also at the same time invites our orthodox Catholic
and Apostolic Church of Christ to unite with the papal throne, thinking
that such union can only be obtained by acknowledging him as supreme
pontiff and the highest spiritual and temporal ruler of the universal
Church, as the only representative of Christ upon earth and the
dispenser of all grace.

III. No doubt every Christian heart ought to be filled with longing for
union of the Churches, and especially the whole orthodox world, being
inspired by a true spirit of piety, according to the divine purpose of
the establishment of the church by the God-man our Savior Christ,
ardently longs for the unity of the Churches in the one rule of faith,
and on the foundation of the apostolic doctrine handed down to us
through the Fathers, \textquotesingle Jesus Christ Himself being the
chief corner stone.\textquotesingle{} {[}1{]} Wherefore she also every
day, in her public prayers to the Lord, prays for the gathering together
of the scattered and for the return of those who have gone astray to the
right way of the truth, which alone leads to the Life of all, the
only-begotten Son and Word of God, our Lord Jesus Christ. {[}2{]}
Agreeably, therefore, to this sacred longing, our orthodox Church of
Christ is always ready to accept any proposal of union, if only the
Bishop of Rome would shake off once for all the whole series of the many
and divers anti-evangelical novelties that have been
\textquotesingle privily brought in\textquotesingle{} to his Church, and
have provoked the sad division of the Churches of the East and West, and
would return to the basis of the seven holy Ecumenical Councils, which,
having been assembled in the Holy Spirit, of representatives of all the
holy Churches of God, for the determination of the right teaching of the
faith against heretics, have a universal and perpetual supremacy in the
Church of Christ. And this, both by her writings and encyclical letters,
the Orthodox Church has never ceased to intimate to the Papal Church,
having clearly and explicitly set forth that so long as the latter
perseveres in her innovations, and the orthodox Church adheres to the
divine and apostolic traditions of Christianity, during which the
Western Churches were of the same mind and were united with the Churches
of the East, so long is it a vain and empty thing to talk of union. For
which cause we have remained silent until now, and have declined to take
into consideration the papal encyclical in question, esteeming it
unprofitable to speak to the ears of those who do not hear. Since,
however, from a certain period the Papal Church, having abandoned the
method of persuasion and discussion, began, to our general astonishment
and perplexity, to lay traps for the conscience of the more simple
orthodox Christians by means of deceitful workers transformed into
apostles of Christ, {[}3{]} sending into the East clerics with the dress
and headcovering of orthodox priests, inventing also divers and other
artful means to obtain her proselytizing objects; for this reason, as in
sacred duty bound, we issue this patriarchal and synodical encyclical,
for a safeguard of the orthodox faith and piety, knowing
\textquotesingle that the observance of the true canons is a duty for
every good man, and much more for those who have been thought worthy by
Providence to direct the affairs of others.\textquotesingle{} {[}4{]}

IV. The union of the separated Churches with herself in one rule of
faith is, as has been said before, a sacred and inward desire of the
holy, catholic and orthodox apostolic Church of Christ; but without such
unity in the faith, the desired union of the Churches becomes
impossible. This being the case, we wonder in truth how Pope Leo XIII,
though he himself also acknowledges this truth, falls into a plain
self-contradiction, declaring, on the one hand, that true union lies in
the unity of faith, and, on the other hand, that every Church, even
after the union, can hold her own dogmatic and canonical definitions,
even when they differ from those of the Papal Church, as the Pope
declares in a previous encyclical, dated November 30, 1894. For there is
an evident contradiction when in one and the same Church one believes
that the Holy Ghost proceeds from the Father, and another that He
proceeds from the Father and the Son; when one sprinkles, and another
baptizes (immerses) thrice in the water; one uses leavened bread in the
sacrament of the Holy Eucharist, and another unleavened; one imparts to
the people of the chalice as well as of the bread, and the other only of
the holy bread; and other things like these. But what this contradiction
signifies, whether respect for the evangelical truths of the holy Church
of Christ and an indirect concession and acknowledgment of them, or
something else, we cannot say.

V. But however that may be, for the practical realization of the pious
longing for the union of the Churches, a common principle and basis must
be settled first of all; and there can be no such safe common principle
and basis other than the teaching of the Gospel and of the seven holy
Ecumenical Councils. Reverting, then, to that teaching which was common
to the Churches of the East and of the West until the separation, we
ought, with a sincere desire to know the truth, to search what the one
holy, catholic and orthodox apostolic Church of Christ, being then
\textquotesingle of the same body,\textquotesingle{} throughout the East
and West believed, and to hold this fact, entire, and unaltered. But
whatsoever has in later times been added or taken away, every one has a
sacred and indispensable duty, if he sincerely seeks for the glory of
God more than for his own glory, that in a spirit of piety he should
correct it, considering that by arrogantly continuing in the perversion
of the truth he is liable to a heavy account before the impartial
judgment-seat of Christ. In saying this we do not at all refer to the
differences regarding the ritual of the sacred services and the hymns,
or the sacred vestments, and the like, which matters, even though they
still vary, as they did of old, do not in the least injure the substance
and unity of the faith; but we refer to those essential differences
which have reference to the divinely transmitted doctrines of the faith,
and the divinely instituted canonical constitution of the administration
of the Churches. \textquotesingle In cases where the thing disregarded
is not the faith (says also the holy Photius), {[}5{]} and is no falling
away from any general and catholic decree, different rites and customs
being observed among different people, a man who knows how to judge
rightly would decide that neither do those who observe them act wrongly,
nor do those who have not received them break the law.\textquotesingle{}
{[}6{]}

VI. And indeed for the holy purpose of union, the Eastern orthodox and
catholic Church of Christ is ready heartily to accept all that which
both the Eastern and Western Churches unanimously professed before the
ninth century, if she has perchance perverted or does not hold it. And
if the Westerns prove from the teaching of the holy Fathers and the
divinely assembled Ecumenical Councils that the then orthodox Roman
Church, which was throughout the West, even before the ninth century
read the Creed with the addition, or used unleavened bread, or accepted
the doctrine of a purgatorial fire, or sprinkling instead of baptism, or
the immaculate conception of the ever-Virgin, or the temporal power, or
the infallibility and absolutism of the Bishop of Rome, we have no more
to say. But if, on the contrary, it is plainly demonstrated, as those of
the Latins themselves, who love the truth, also acknowledge, that the
Eastern and orthodox catholic Church of Christ holds fast the anciently
transmitted doctrines which were at that time professed in common both
in the East and the West, and that the Western Church perverted them by
divers innovations, then it is clear, even to children, that the more
natural way to union is the return of the Western Church to the ancient
doctrinal and administrative condition of things; for the faith does not
change in any way with time or circumstances, but remains the same
always and everywhere, for \textquotesingle there is one body and one
Spirit,\textquotesingle{} it is said, \textquotesingle even as ye are
called in one hope of your calling; one Lord, one faith, one baptism,
one God and Father of all, who is above all, and through all, and in you
all." {[}7{]}

VII. So then the one holy, catholic and apostolic Church of the seven
Ecumenical Councils believed and taught in accordance with the words of
the Gospel that the Holy Ghost proceeds from the Father; but in the
West, even from the ninth century, the holy Creed, which was composed
and sanctioned by Ecumenical Councils, began to be falsified, and the
idea that the Holy Ghost proceeds \textquotesingle also from the
Son\textquotesingle{} to be arbitrarily promulgated. And certainly Pope
Leo XIII is not ignorant that his orthodox predecessor and namesake, the
defender of orthodoxy, Leo III, in the year 809 denounced synodically
this anti-evangelical and utterly lawless addition, \textquotesingle and
from the Son\textquotesingle{} \emph{(filioque);} and engraved on two
silver plates, in Greek and Latin, the holy Creed of the first and
second Ecumenical Councils, entire and without any addition; having
written moreover, \textquotesingle These words I, Leo, have set down for
love and as a safeguard of the orthodox faith\textquotesingle{}
(\emph{Haec Leo posui amore et cautela fidei
orthodoxa}\textquotesingle). {[}8{]}

Likewise he is by no means ignorant that during the tenth century, or at
the beginning of the eleventh, this anti-evangelical and lawless
addition was with difficulty inserted officially into the holy Creed at
Rome also, and that consequently the Roman Church, in insisting on her
innovations, and not coming back to the dogma of the Ecumenical
Councils, renders herself fully responsible before the one holy,
catholic and apostolic Church of Christ, which holds fast that which has
been received from the Fathers, and keeps the deposit of the faith which
was delivered to it unadulterated in all things, in obedience to the
Apostolic injunction: \textquotesingle That good thing which was
committed unto thee keep by the Holy Ghost which dwelleth in
us\textquotesingle; \textquotesingle avoiding profane and vain
babblings, and oppositions of science falsely so called: which some
professing have erred concerning the faith." {[}9{]}

VIII. The one holy, catholic and apostolic Church of the first seven
Ecumenical Councils baptized by three immersions in the water, and the
Pope Pelagius speaks of the triple immersion as a command of the Lord,
and in the thirteenth century baptism by immersions still prevailed in
the West; and the sacred fonts themselves, preserved in the more ancient
churches in Italy, are eloquent witnesses on this point; but in later
times sprinkling or effusion, being privily brought in, came to be
accepted by the Papal Church, which still holds fast the innovation,
thus also widening the gulf which she has opened; but we Orthodox,
remaining faithful to the apostolic tradition and the practice of the
seven Ecumenical Councils, \textquotesingle stand fast, contending for
the common profession, the paternal treasure of the sound
faith.\textquotesingle{} {[}10{]}

IX. The one holy, catholic and apostolic Church of the seven Ecumenical
Councils, according to the example of our Savior, celebrated the divine
Eucharist for more than a thousand years throughout the East and West
with leavened bread, as the truth-loving papal theologians themselves
also bear witness; but the Papal Church from the eleventh century made
an innovation also in the sacrament of the divine Eucharist by
introducing unleavened bread.

X. The one holy, catholic and apostolic Church of the seven Ecumenical
Councils held that the precious gifts are consecrated after the prayer
of the invocation of the Holy Ghost by the blessing of the priest, as
the ancient rituals of Rome and Gaul testify; nevertheless afterwards
the Papal Church made an innovation in this also, by arbitrarily
accepting the consecration of the precious gifts as taking place along
with the utterance of the Lord\textquotesingle s words:
\textquotesingle Take, eat; this is my body\textquotesingle: and
\textquotesingle Drink ye all of it; for this is my
blood.\textquotesingle{} {[}11{]}

XI. The one holy, catholic and apostolic Church of the seven Ecumenical
Councils, following the Lord\textquotesingle s command,
\textquotesingle Drink ye all of it,\textquotesingle{} {[}12{]} imparted
also of the holy chalice to all; but the Papal Church from the ninth
century downwards has made an innovation in this rite also, by depriving
the laity of the holy chalice, contrary to the Lord\textquotesingle s
command and the universal practice of the ancient Church, as well as the
express prohibition of many ancient orthodox bishops of Rome.

XII. The one holy, catholic and apostolic Church of the seven Ecumenical
Councils, walking according to the divinely inspired teaching of the
Holy Scripture and the old apostolic tradition, prays and invokes the
mercy of God for the forgiveness and rest of those
\textquotesingle which have fallen asleep in the Lord\textquotesingle;
{[}13{]} but the Papal Church from the twelfth century downwards has
invented and heaped together in the person of the Pope, as one
singularly privileged, a multitude of innovations concerning purgatorial
fire, a superabundance of the virtues of the saints, and the
distribution of them to those who need them, and the like, setting forth
also a full reward for the just before the universal resurrection and
judgment.

XIII. The one holy, catholic and apostolic Church of the seven
Ecumenical Councils teaches that the supernatural incarnation of the
only-begotten Son and Word of God, of the Holy Ghost and the Virgin
Mary, is \emph{alone} pure and immaculate; but the Papal Church scarcely
forty years ago again made an innovation by laying down a novel dogma
concerning the immaculate conception of the Mother of God and
ever-Virgin Mary, which was unknown to the ancient Church (and strongly
opposed at different times even by the more distinguished among the
papal theologians).

XIV. Passing over, then, these serious and substantial differences
between the two churches respecting the faith, which differences, as has
been said before, were created in the West, the Pope in his encyclical
represents the question of the primacy of the Roman Pontiff as the
principal and, so to speak, only cause of the dissension, and sends us
to the sources, that we may make diligent search as to what our
forefathers believed and what the first age of Christianity delivered to
us. But having recourse to the fathers and the Ecumenical Councils of
the Church of the first nine centuries, we are fully persuaded that the
Bishop of Rome was never considered as the supreme authority and
infallible head of the Church, and that every bishop is head and
president of his own particular Church, subject only to the synodical
ordinances and decisions of the Church universal as being alone
infallible, the Bishop of Rome being in no wise excepted from this rule,
as Church history shows. Our Lord Jesus Christ alone is the eternal
Prince and immortal Head of the Church, for \textquotesingle He is the
Head of the body, the Church," {[}14{]} who said also to His divine
disciples and apostles at His ascension into heaven,
\textquotesingle Lo, I am with you alway, even unto the end of the
world.\textquotesingle{} {[}15{]} In the Holy Scripture the Apostle
Peter, whom the Papists, relying on apocryphal books of the second
century, the pseudo-Clementines, imagine with a purpose to be the
founder of the Roman Church and their first bishop, discusses matters as
an equal among equals in the apostolic synod of Jerusalem, and at
another time is sharply rebuked by the Apostle Paul, as is evident from
the Epistle to the Galatians. {[}16{]} Moreover, the Papists themselves
know well that the very passage of the Gospel to which the Pontiff
refers, \textquotesingle Thou art Peter, and upon this rock I will build
my Church,\textquotesingle{} {[}17{]} is in the first centuries of the
Church interpreted quite differently, in a spirit of orthodoxy, both by
tradition and by all the divine and sacred Fathers without exception;
the fundamental and unshaken rock upon which the Lord has built His own
Church, against which the gates of hell shall not prevail, being
understood metaphorically of Peter\textquotesingle s true confession
concerning the Lord, that \textquotesingle He is Christ, the Son of the
living God.\textquotesingle{} {[}18{]} Upon this confession and faith
the saving preaching of the Gospel by all the apostles and their
successors rests unshaken. Whence also the Apostle Paul, who had been
caught up into heaven, evidently interpreting this divine passage,
declares the divine inspiration, saying: \textquotesingle According to
the grace of God which is given unto me, as a wise master-builder, I
have laid the foundation, and another buildeth thereon. For other
foundation can no man lay than that is laid, which is Jesus
Christ.\textquotesingle{} {[}19{]} But it is in another sense that Paul
calls all the apostles and prophets together the foundation of the
building up in Christ of the faithful; that is to say, the members of
the body of Christ, which is the Church; {[}20{]} when he writes to the
Ephesians: \textquotesingle Now therefore ye are no more strangers and
foreigners, but fellow-citizens with the saints and of the house hold of
God; and are built upon the foundation of the apostles and prophets,
Jesus Christ Himself being the chief corner stone.\textquotesingle{}
{[}21{]} Such, then, being the divinely inspired teaching of the
apostles respecting the foundation and Prince of the Church of God, of
course the sacred Fathers, who held firmly to the apostolic traditions,
could not have or conceive any idea of an absolute primacy of the
Apostle Peter and the bishops of Rome; nor could they give any other
interpretation, totally unknown to the Church, to that passage of the
Gospel, but that which was true and right; nor could they arbitrarily
and by themselves invent a novel doctrine respecting excessive
privileges of the Bishop of Rome as successor, if so be, of Peter;
especially whilst the Church of Rome was chiefly founded, not by Peter,
whose apostolic action at Rome is totally unknown to history, but by the
heaven-caught apostle of the Gentiles, Paul, through his disciples,
whose apostolic ministry in Rome is well known to all. {[}22{]}

XV. The divine Fathers, honoring the Bishop of Rome only as the bishop
of the capital city of the Empire, gave him the honorary prerogative of
presidency, considering him simply as the bishop first in order, that
is, first among equals; which prerogative they also assigned afterwards
to the Bishop of Constantinople, when that city became the capital of
the Roman Empire, as the twenty-eighth canon of the fourth Ecumenical
Council of Chalcedon bears witness, saying, among other things, as
follows: \textquotesingle We do also determine and decree the same
things respecting the prerogatives of the most holy Church of the said
Constantinople, which is New Rome. For the Fathers have rightly given
the prerogative to the throne of the elder Rome, because that was the
imperial city. And the hundred and fifty most religious bishops, moved
by the same consideration, assigned an equal prerogative to the most
holy throne of New Rome.\textquotesingle{} From this canon it is very
evident that the Bishop of Rome is equal in honor to the Bishop of the
Church of Constantinople and to those other Churches, and there is no
hint given in any canon or by any of the Fathers that the Bishop of Rome
alone has ever been prince of the universal Church and the infallible
judge of the bishops of the other independent and self-governing
Churches, or the successor of the Apostle Peter and vicar of Jesus
Christ on earth.

XVI. Each particular self-governing Church, both in the East and West,
was totally independent and self-administered in the time of the Seven
Ecumenical Councils. And just as the bishops of the self-governing
Churches of the East, so also those of Africa, Spain, Gaul, Germany and
Britain managed the affairs of their own Churches, each by their local
synods, the Bishop of Rome having no right to interfere, and he himself
also was equally subject and obedient to the decrees of synods. But on
important questions which needed the sanction of the universal Church an
appeal was made to an Ecumenical Council, which alone was and is the
supreme tribunal in the universal Church. Such was the ancient
constitution of the Church; but the bishops were independent of each
other and each entirely free within his own bounds, obeying only the
syndical decrees, and they sat as equal one to another in synods.
Moreover, none of them ever laid claim to monarchical rights over the
universal Church; and ii sometimes certain ambitious bishops of Rome
raised excessive claims to an absolutism unknown to the Church, such
were duly reproved and rebuked The assertion therefore of Leo XIII, when
he says in his Encyclical that before the period of the great Photius
the name of the Roman throne was holy among all the peoples of the
Christian world, and that the East, like the West, with one accord and
without opposition, was subject to the Roman pontiff as lawful
successor, so to say, of the Apostle Peter, and consequently vicar of
Jesus Christ on earth is proved to be inaccurate and a manifest error.

XVII. During the nine centuries of the Ecumenical Councils the Eastern
Orthodox Church never recognized the excessive claims of primacy on the
part of the bishops of Rome, nor consequently did she ever submit
herself to them, as Church history plainly bears witness. The
independent relation of the East to the West is clearly and manifestly
shown also by those few and most significant words of Basil the Great,
which he writes in a letter to the holy Eusebius, Bishop of Samosata:
\textquotesingle For when haughty characters are courted, it is their
nature to become still more disdainful. For if the Lord be merciful to
us, what other assistance do we need? But if the wrath of God abide on
us, what help is there for us from Western superciliousness? Men who
neither know the truth nor can bear to learn it, but being prejudiced by
false suspicions, they act now as they did before in the case of
Marcellus.\textquotesingle{} {[}23{]} The celebrated Photius, therefore,
the sacred Prelate and luminary of Constantinople, defending this
independence of the Church of Constantinople after the middle of the
ninth century, and foreseeing the impending perversion of the
ecclesiastical constitution in the West, and its defection from the
orthodox East, at first endeavored in a peaceful manner to avert the
danger; but the Bishop of Rome, Nicholas 1, by his uncanonical
interference with the East, beyond the bounds of his diocese, and by the
attempt which he made to subdue the Church of Constantinople to himself,
pushed maners to the verge of the grievous separation of the Churches.
The first seeds of these claims of a papal absolutism were scattered
abroad in the pseudo-Clementines, and were cultivated, exactly at the
epoch of this Nicholas, in the so-called \emph{pseudo-lsidorian
decrees,} which are a farrago of spurious and forged royal decrees and
letters of ancient bishops of Rome, by which, contrary to the truth of
history and the established constitution of the Church, it was purposely
promulgated that, as they said, Christian antiquity assigned to the
bishops of Rome an unbounded authority over the universal Church.

XVIII. These facts we recall with sorrow of heart, inasmuch as the Papal
Church, though she now acknowledges the spuriousness and forged
character of those decrees on which her excessive claims are grounded,
not only stubbornly refuses to come back to the canons and decrees of
the Ecumenical Councils, but even in the expiring years of the
nineteenth century has widened the existing gulf by officially
proclaiming, to the astonishment of the Christian world, that the Bishop
of Rome is even infallible. The orthodox Eastern and catholic Church of
Christ, with the exception of the Son and Word of God, who was ineffably
made man, knows no one infallible upon earth. Even the Apostle Peter
himself, whose successor the Pope thinks himself to be, thrice denied
the Lord, and was twice rebuked by the Apostle Paul, as not walking
uprightly according to the truth of the Gospel. {[}24{]} Afterwards the
Pope Liberius, in the fourth century, subscribed an Arian confession;
and likewise Zosimus, in the fifth century, approved an heretical
confession, denying original sin. Virgilius, in the sixth century, was
condemned for wrong opinions by the fifth Council; and Honorius, having
fallen into the Monothelite heresy, was condemned in the seventh century
by the sixth Ecumenical Council as a heretic, and the popes who
succeeded him acknowledged and accepted his condemnation.

XIX. With these and such facts in view, the peoples of the West,
becoming gradually civilized by the diffusion of letters, began to
protest against innovations, and to demand (as was done in the fifteenth
century at the Councils of Constance and Basle) the return to the
ecclesiastical constitution of the first centuries, to which, by the
grace of God, the orthodox Churches throughout the East and North, which
alone now form the one holy, catholic and apostolic Church of Christ,
the pillar and ground of the truth, remain, and will always remain,
faithful. The same was done in the seventeenth century by the learned
Gallican theologians, and in the eighteenth by the bishops of Germany;
and in this present century of science and criticism, the Christian
conscience rose up in one body in the year 1870, in the persons of the
celebrated clerics and theologians of Germany, on account of the novel
dogma of the infallibility of the Popes, issued by the Vatican Council,
a consequence of which rising is seen in the formation of the separate
religious communities of the old Catholics, who, having disowned the
papacy, are quite independent of it.

XX. In vain, therefore, does the Bishop of Rome send us to the sources
that we may seek diligently for what our forefathers believed and what
the first period of Christianity delivered to us. In these sources we,
the orthodox, find the old and divinely-transmitted doctrines, to which
we carefully hold fast to the present time, and nowhere do we find the
innovations which later times of empty mindedness brought forth in the
West, and which the Papal Church having adopted retains till this very
day. The orthodox Eastern Church then justly glories in Christ as being
the Church of the seven Ecumenical Councils and of the first nine
centuries of Christianity, and therefore the one holy, catholic and
apostolic Church of Christ, \textquotesingle the pillar and ground of
the truth\textquotesingle; {[}25{]} but the present Roman Church is the
Church of innovations, of the falsification of the writings of the
Church Fathers, and of the misinterpretation of the Holy Scripture and
of the decrees of the holy councils, for which she has reasonably and
justly been disowned, and is still disowned, so far as she remains in
her error. \textquotesingle For better is a praiseworthy war than a
peace which separates from God,\textquotesingle{} as Gregory of
Nazianzus also says.

XXI. Such are, briefly, the serious and arbitrary innovations concerning
the faith and the administrative constitution of the Church, which the
Papal Church has introduced and which, it is evident, the Papal
Encyclical purposely passes over in silence. These innovations, which
have reference to essential points of the faith and of the
administrative system of the Church, and which are manifestly opposed to
the ecclesiastical condition of the first nine centuries, make the
longed-for union of the Churches impossible: and every pious and
orthodox heart is filled with inexpressible sorrow on seeing the Papal
Church disdainfully persisting in them, and not in the least
contributing to the sacred purpose of union by rejecting those heretical
innovations and coming back to the ancient condition of the one holy,
catholic and apostolic Church of Christ, of which she also at that time
formed a part.

XXII. But what are we to say of all that the Roman Pontiff writes when
he addresses the glorious Slavonic nations? No one, indeed, has ever
denied that by the virtue and the apostolic toils of SS. Cyril and
Methodius the grace of salvation was vouchsafed to not a few of the
Slavonic peoples: but history testifies that at the period of the great
Photius those Greek apostles to the Slavs and intimate friends of that
divine Father, setting out from Thessalonica, were sent to convert the
Slavonic tribes not from Rome but from Constantinople, where moreover
they had been trained, living as monks in the monastery of St.
Polychronius. It is therefore utterly incoherent which is proclaimed in
the Roman Pontiff\textquotesingle s Encyclical, that, as he says, a
kindly relation and mutual sympathy was brought about between the
Slavonic tribes and the pontiffs of the Roman Church; for even if the
Pope is ignorant of it, history nevertheless explicitly proclaims that
these sacred apostles to the Slavs of whom we speak, encountered greater
difficulties in their work from the bishops of Rome through their
excommunications and opposition, and were more cruelly persecuted by the
Frankish papal bishops than by the heathen inhabitants of those
countries. Certainly the Pope knows well that the blessed Methodius
having departed to the Lord, two hundred of the most distinguished of
his disciples\textquotesingle{} after many struggles against the
opposition of the Roman Pontiffs, were driven out of Moravia and led
away by military force beyond its boundaries, from whence afterwards
they were dispersed into Bulgaria and elsewhere. And he knows also that
with the expulsion of the more erudite Slavonic clergy, the ritual of
the East, as well as the Slavonic language then in use, were also driven
out, and in process of time all vestige of orthodoxy was effaced from
those provinces, and all these things done with the official cooperation
of the bishops of Rome m a manner not the least honorable to the
holiness of the episcopal dignity. But notwithstanding all this
despiteful treatment, the orthodox Slavonic Churches, the beloved
daughters of the orthodox East, and especially the great and glorious
Church of divinely preserved Russia, having been preserved harmless by
the grace of God, have kept, and will keep till the end of the ages, the
orthodox faith, and stand forth conspicuous testimonies of the liberty
that is in Christ. In vain, therefore, does the Papal Encyclical promise
to the Slavonic Churches prosperity and greatness, because by the
goodwill of the most gracious God they already possess these blessings,
and such as these, standing firm m the orthodoxy of their fathers and
glorifying in it in Christ.

XXIII. These things being so, and being indisputably proved by
ecclesiastical history, we, anxious as it is our duty to be, address
ourselves to the peoples of the West, who through ignorance of the true
and impartial history of ecclesiastical matters, being credulously led
away, follow the anti-evangelical and utterly lawless innovations of the
papacy, having been separated and continuing far from the one holy,
catholic and apostolic orthodox Church of Christ, which is
\textquotesingle the Church of the living God, the pillar and ground of
the truth, {[}26{]} in which also their gracious ancestors and
forefathers shone by their piety and orthodoxy of faith, having been
faithful and precious members of it during nine whole centuries,
obediently following and walking according to the decrees of the
divinely assembled Ecumenical Councils.

XXIV. Christ-loving peoples of the glorious countries of the West! We
rejoice on the one hand seeing that you have a zeal for Christ, being
led by this right persuasion, \textquotesingle that without faith in
Christ it is impossible to please God\textquotesingle; {[}27{]} but on
the other hand it is self-evident to every right-thinking person that
the salutary faith in Christ ought by all means to be right in
everything, and in agreement with the Holy Scripture and the apostolic
traditions, upon which the teaching of the divine Fathers and the seven
holy, divinely assembled Ecumenical Councils is based. It is moreover
manifest that the universal Church of God, which holds fast in its bosom
unique unadulterated and entire this salutary faith as a divine deposit,
just as it was of old delivered and unfolded by the God-bearing Fathers
moved by the Spirit, and formulated by them during the first nine
centuries, is one and the same for ever, and not manifold and varying
with the process of time: because the gospel truths are never
susceptible to alteration or progress in course of time, like the
various philosophical systems; \textquotesingle for Jesus Christ is the
same yesterday, and to day, and for ever.\textquotesingle{} {[}28{]}
Wherefore also the holy Vincent, who was brought up on the milk of the
piety received from the fathers in the monastery of Lérins in Gaul, and
flourished about the middle of the fifth century, with great wisdom and
orthodoxy characterizes the true catholicity of the faith and of the
Church, saying: \textquotesingle In the catholic Church we must
especially take heed to hold that which has been believed everywhere at
all times, and by all. For this is truly and properly catholic, as the
very force and meaning of the word signifies, which moreover comprehends
almost everything universally. And that we shall do, if we walk
following universality, antiquity, and consent.\textquotesingle{}
{[}29{]} But, as has been said before, the Western Church, from the
tenth century downwards, has privily brought into herself through the
papacy various and strange and heretical doctrines and innovations, and
so she has been torn away and removed far from the true and orthodox
Church of Christ. How necessary, then, it is for you to come back and
return to the ancient and unadulterated doctrines of the Church in order
to attain the salvation in Christ after which you press, you can easily
understand if you intelligently consider the command of the
heaven-ascended Apostle Paul to the Thessalonians, saying:
\textquotesingle Therefore, brethren, stand fast, and hold the
traditions which ye have been taught, whether by word, or our
epistle\textquotesingle; {[}30{]} and also what the same divine apostle
writes to the Galatians saying: \textquotesingle I marvel that ye are so
soon removed from him that called you into the grace of Christ unto
another gospel: which is not another; but there be some that trouble
you, and would pervert the gospel of Christ.\textquotesingle{} {[}31{]}
But avoid such perverters of the evangelical truth, \textquotesingle For
they that are such serve not our Lord Jesus Christ, but their own belly;
and by good words and fair speeches deceive the hearts of the
simple;{[}32{]} and come back for the future into the bosom of the holy,
catholic and apostolic Church of God, which consists of all the
particular holy Churches of God, which being divinely planted, like
luxuriant vines throughout the orthodox world, are inseparably united to
each other in the unity of the one saving faith in Christ, and in the
bond of peace and of the Spirit, that you may obtain the
highly-to-be-praised and most glorious name of our Lord and God and
Savior Jesus Christ, who suffered for the salvation of the world, may be
glorified among you also.

XXV. But let us, who by the grace and goodwill of the most gracious God
are precious members of the body of Christ, that is to say of His one
holy, catholic and apostolic Church, hold fast to the piety of our
fathers, handed down to us from the apostles. Let us all beware of false
apostles, who, coming to us in sheep\textquotesingle s clothing, attempt
to entice the more simple among us by various deceptive promises,
regarding all things as lawful and allowing them for the sake of union,
provided only that the Pope of Rome be recognized as supreme and
infallible ruler and absolute sovereign of the universal Church, and
only representative of Christ on earth, and the source of all grace. And
especially let us, who by the grace and mercy of God have been appointed
bishops, pastors, and teachers of the holy Churches of God,
\textquotesingle take heed unto ourselves,---and to all the flock, over
which the Holy Ghost hath made us overseers, to feed the Church of God,
which He hath purchased with His own blood,\textquotesingle{} {[}33{]}
as they that must give account. \textquotesingle Wherefore let us
comfort ourselves together, and edify one another.\textquotesingle{}
{[}34{]} \textquotesingle And the God of all grace, who hath called us
unto His eternal glory by Christ Jesus ... make us perfect, stablish,
strengthen, settle us,\textquotesingle{} {[}35{]} and grant that all
those who are without and far away from the one holy, catholic and
orthodox fold of His reasonable sheep may be enlightened with the light
of His grace and the acknowledging of the truth. To Him be glory and
dominion for ever and ever.

Amen.

In the Patriarchal Palace of Constantinople, in the month of August of
the year of grace MDCCCXCV.

+ ANTHIMOS of Constantinople, beloved brother and intercessor in Christ
our God.

+ NICODEMOS of Cyzicos, beloved brother and intercessor in Christ our
God.

+ PHILOTHEOS of Nicomedia, beloved brother and intercessor in Christ our
God.

+ JEROME of Nicea, beloved brother and intercessor in Christ our God.

+ NATHANAEL of Prusa, beloved brother and intercessor of Christ our God.

+ BASIL of Smyrna, beloved brother and intercessor in Christ our God.

+ STEPHEN of Philadelphia, beloved brother and intercessor in Christ our
God.

+ ATHANASIOS of Lemnos, beloved brother and intercessor in Christ our
God.

+ BESSARION of Dyrrachium, beloved brother and intercessor in Christ our
God.

+ DOROTHEOS of Belgrade, beloved brother and intercessor in Christ our
God.

+ NICODEMOS of Elasson, beloved brother and intercessor in Christ our
God.

+ SOPHRONIOS of Carpathos and Cassos, beloved brother and intercessor in
Christ our God.

+ DIONYSIOS of Eleutheropolis, beloved brother and intercessor in Christ
our God.

\section{Endnotes}\label{endnotes}

1. Eph. 2:20.

2. John 14:6.

3. II Cor. 11:13.

4. Phot. \emph{Epist.} iii. 10.

5. Patriarch of Constantinople; c. 800.

6. Phot. \emph{Epist} iii. 6.

7. Eph. 4:5-6.

8. See life of Leo 111 by Athanasius, presbyter and librarian at Rome,
in his \emph{Lives of the Popes.} The holy Photius also, making mention
of this invective of the orthodox Pope of Rome, Leo III, against the
holders of the erroneous doctrine, in his renowned letter to the
Metropolitan of Acquileia, expresses himself as follows:
\textquotesingle For (not to mention those who were before him) Leo the
elder, prelate of Rome, as well as Leo the younger after him, shew
themselves to be of the same mind with the catholic and apostolic
Church, with the holy prelates their predecessors, and with the
apostolic commands; the one having contributed much to the assembling of
the fourth holy Ecumenical Council, both by the sacred men who were sent
to represent him, and by his letter, through which both Nestorius and
Eutyches were overthrown; by which letter he moreover, in accordance
with previous synodical decrees, declared the Holy Ghost to proceed from
the Father, but not also "from the Son." And in like manner Leo the
younger, his counterpart in faith as well as in name. This latter
indeed, who was ardently zealous for true piety, in order that the
unspotted pattern of true piety might not in any way whatever be
falsified by a barbarous language, published it in Greek, as has already
been said in the beginning, to the people of the West, that they might
thereby glorify and preach aright the Holy Trinity. And not only by word
and command, but also, having inscribed and exposed it to the sight of
all on certain shields specially made, as on certain monuments, he fixed
it at the gates of the Church, in order that every person might easily
learn the uncontaminated faith, and in order that no chance whatever
might be left to secret forgers and innovators of adulterating the piety
of us Christians, and of bringing in the Son besides the Father as a
second cause of the Holy Spirit, who proceeds from the Father with honor
equal to that of the begotten Son. And it was not these two holy men
alone, who shone brightly in the West, who preserved the faith free from
innovation; for the Church is not in such want as that of Western
preachers; but there is also a host of them not easily counted who did
likewise.\textquotesingle---\emph{Epist. v. 53.}

\emph{9.} III Tim. 1:14; 1 Tim. 6:20-21.

10. St. Basil the Great, \emph{Ep.} 243, To the Bishops of Italy and
Gaul.

11. Matt. 26:26, 28

12. Matt. 26:28.

13. Matt. 26:31; Heb. 11:39-40; II Tim. 4:8; II Macc. 12:45.

14. Col. 1:18.

15. Matt. 28:20.

16. Gal. 2:11.

17. Matt. 16:18.

18. Matt. 16:16.

19. 1 Cor. 3:10, 11.

20. Col. 1:24.

21. Eph. 2:19, 20. Cp. 1 Pet. 2:4; Rev. 21:14.

22. See Acts of the Apostles 28:15, Rom. 15:15-16; Phil. 1:13.

23. \emph{Epist. 239.}

24. Gal. 2:11.

25. I Tim. 3:15.

26. I Tim. 3:15.

27. Heb. 11:6.

28. Heb. 13:8.

29. \textquotesingle{}\emph{In ipsa item Catholica Ecclesia magnopere
curandum est, ut teneamus, quod ubique quod semper ab omnibus creditum
est. Hoc est enim vere proprieque Catholicum (quod ipsa vis nominis
ratioque declarat), quod omnia fere universaliter comprehendit. Sed hoc
fiet si sequimur universalitatem, antiquitatem,
consensionem\textquotesingle{}} (Vincentii Lirinensis
\emph{Commonitorium pro CatholicEe fidei antiquitate et universalitate}
cap. iii, cf. cap. viii and xiv).

30. 1Thess.2:15.

31. Gal. 1:6-7.

32. Rom. 16:18.

33. Acts 20:28.

34. I Thess. 5:11.

35. I Pet. 5:10.
